Representation engineering shows that large language models encode sycophancy as a linear direction in activation space. We extract a contrastive sycophancy direction ($\dsyc$) from style-controlled pairs on two model families, achieving cross-validated AUROC $0.827$ on Qwen3-8B and $0.779$ on Llama-3-8B. We then test three $\dsyc$-based reward formulations with Group Relative Policy Optimization (GRPO) on Qwen3-8B. Every formulation fails. The cosine reward cannot separate valid revision from sycophantic capitulation in the narrow online activation gap. The probe reward creates a stubbornness trap: the model learns that never revising its answer maximizes reward, and the true-positive flip rate drops to zero within 500 steps. When we remove $\dsyc$ entirely and train with a purely behavioral reward, only 4 of 25 checkpoints across two seeds exceed vanilla, versus 17/17 for supervised fine-tuning ($p < 0.0001$). Only supervised fine-tuning with dense per-token supervision succeeds. The root cause is a semantic gap between discriminability and reward utility: $\dsyc$ encodes \emph{whether} belief changed, not \emph{whether it should have}. Valid revision and sycophantic capitulation produce overlapping activations because both involve the model shifting its stated position toward the user's. A single scalar reward per completion, whether from representations or behavioral classifiers, cannot convey which belief changes are warranted. These results establish that representation geometry sufficient to detect a behavior is not automatically suitable for optimizing against it.
